% D.2.47: complete draft held for independent proof/code audit.
\paragraph{Theorem D.2.47 (Proved: the general triangular-cusp component intersection form).}
Let $N_L$ be the smooth triangular Mumford cap of D.2.43,
with period sublattice $L\subset\mathbb Z^2$ of finite index
$D>0$. Write $V=\mathbb Z^2/L$ for its vertices and
$E_i$, $i\in V$, for its reduced irreducible central divisors.
Their normalizations $\nu_i:P_i\to E_i\subset N_L$ are the
degree-six toric del Pezzo surface, even when a divisor has
self-gluings. Use the cyclic ray order
\[
 (r_0,\ldots,r_5)=
       (e_1,e_2,e_2-e_1,-e_1,-e_2,e_1-e_2).                 \tag{D47.1}
\]
For a coefficient vector $u=(u_i)_{i\in V}$ put
$E(u)=\sum_i u_iE_i$ and
$a_{i,r}(u)=u_{i+r}-u_i$. The entire integral trilinear
form on these compact divisor classes is
\[
\begin{split}
 I_L(u,v,w)=\sum_{i\in V}u_i\Bigl(
 &-\sum_{r=0}^5 a_{i,r}(v)a_{i,r}(w)\\
 &+\sum_{r=0}^5\bigl(a_{i,r}(v)a_{i,r+1}(w)
                    +a_{i,r+1}(v)a_{i,r}(w)\bigr)\Bigr),
\end{split}                                                   \tag{D47.2}
\]
where ray indices are cyclic. It is symmetric in all three
arguments, and the constant vector annihilates it.
Equivalently, the cubic is
\[
\begin{split}
 I_L(u,u,u)={}&6\sum_{i\in V}u_i^3
       -3\sum_{\{i,j\}\in\mathcal E}
                      (u_i^2u_j+u_iu_j^2)\\
       &+6\sum_{(i,j,k)\in\mathcal T}u_iu_ju_k.
\end{split}                                                   \tag{D47.3}
\]
Here $\mathcal E$ is the set of the $3D$ quotient edge
\emph{cells}, indexed by $i$ and the first three directions
in (D47.1), and $\mathcal T$ consists of the $2D$ triangle
cells. Loops, multiple edges and repeated vertex labels in
a triangle are retained with their actual cell multiplicities.
They are not replaced by a simple graph or a simplicial
complex with distinct labels.
The Chern and Pontryagin evaluations on every component are
\[
 \left\langle c_2(TN_L),\nu_{i*}[P_i]\right\rangle=0,
 \qquad
 \left\langle p_1(TN_L),\nu_{i*}[P_i]\right\rangle=0.         \tag{D47.4}
\]
All these formulas remain the same when the cap is included
in a closed smooth complex threefold.

For $L=\mathbb Ze_1\oplus d\mathbb Ze_2$, the entire
component trilinear form is zero for every $d>0$.
It is not zero for arbitrary period lattices: for
$L=2\mathbb Z^2$ and its four vertices,
\[
 E_i^3=6,\qquad E_i^2E_j=-2\ (i\ne j),\qquad
 E_iE_jE_k=2\ (i,j,k\text{ distinct}).                      \tag{D47.5}
\]
This is the component part of the cohomology ring and
the indicated characteristic pairings. It is not asserted
to determine torsion products, the full fundamental-class
invariant, or a global diffeomorphism type.

\emph{Proof: Cartier restriction with self-gluings.}
Denote the six toric boundary curves of $P_i$ by $C_{i,r}$.
The fan (D47.1) is the fan of $\mathbb P^2$ subdivided
at its three torus-fixed points. Consecutive rays have
determinant one and $r_{j-1}+r_{j+1}=r_j$. Consequently
\[
 C_{i,r}^2=-1,\qquad
 C_{i,r}C_{i,r+1}=1,\qquad
 C_{i,r}C_{i,s}=0\quad(s\not\equiv r,r\pm1).               \tag{D47.6}
\]
These also follow directly from the monomial coordinates
on the six smooth affine toric charts.

For $j\ne i$, the pullback of the Cartier divisor $E_j$
to $P_i$ is precisely the sum of those boundary curves
whose adjacent label is $j$. The total reduced central
divisor is $\sum_jE_j=\operatorname{div}(q)$, so its
absolute line bundle on the cap is trivial. Restricting
this identity determines the remaining pullback and gives
the equality of line bundles, or their divisor classes,
\[
 \nu_i^*\mathcal O(E(u))
       =\mathcal O_{P_i}\left(
                   \sum_{r=0}^5(u_{i+r}-u_i)C_{i,r}\right).
                                                               \tag{D47.7}
\]
When an edge returns to the same component, its coefficient
is zero. In local normal-crossing coordinates, the same
irreducible Cartier divisor then contains both branches;
the second branch contributes exactly the correction to
the normal line of the first branch. Thus (D47.7) does
not apply embedded-divisor adjunction incorrectly to a
self-glued component.

The fundamental cycle of $E_i$ is $\nu_{i*}[P_i]$ with
degree one. Intersection with the compact divisor $E(u)$,
then the projection formula, therefore gives
\[
 E(u)E(v)E(w)=\sum_i u_i
           \int_{P_i}c_1\nu_i^*\mathcal O(E(v))
                       c_1\nu_i^*\mathcal O(E(w)).          \tag{D47.8}
\]
Equations (D47.6)--(D47.7) give (D47.2) integrally.
It is symmetric because these are products of actual
Cartier classes on the smooth threefold. A constant
coefficient vector has trivial restriction on every
$P_i$, so annihilates the form.

There is a relative-cohomology distinction here. The
divisors have compact support and define classes in
$H^2(N_L,\partial N_L;\mathbb Z)$; their triple product
is evaluated in relative degree six. Although
$\sum_iE_i$ has a trivial absolute line bundle on the
cap, its compact-support class need not be zero: the
boundary value of $q$ winds around the origin.
The assertion is that the constant vector annihilates
the component intersection form, not that the sum of
the component fundamental cycles vanishes in homology.
For an inclusion into a closed threefold, the same
compact cycles and local divisor restrictions compute
the same products.

\emph{The quotient cell formula and examples.}
Set $v=w=u$ in (D47.2) and expand. A vertex contributes
$6u_i^3$. Pairing the two directed incidences of every
edge gives $-3(u_i^2u_j+u_iu_j^2)$. Each triangle has
three corner incidences and contributes $6u_iu_ju_k$.
This proves (D47.3), with all incidence multiplicities
retained, even when labels coincide.
One explicit set of triangle cells is
\[
 (i,i+e_1,i+e_2),\qquad
 (i+e_1+e_2,i+e_1,i+e_2)\quad(i\in V).                    \tag{D47.9}
\]
The integral trilinear form is already specified by
(D47.2); no divisibility assumption based on a rational
polarization is needed.

If $e_1\in L$ and $V$ is the cyclic sequence of length
$d$ in the $e_2$ direction, the $e_1$ loop terms cancel
the vertex terms. The other two edge families have the
same adjacent labels and cancel the two triangle
families in (D47.9). Hence the cubic, and by (D47.2)
the whole trilinear form, vanish for every $d$, including
$d=1,2$. For $L=2\mathbb Z^2$, each vertex sees the
other three labels twice in its hexagon. Substitution
in (D47.2) gives (D47.5), so a vanishing conclusion
cannot be extrapolated from the cyclic period case.

\emph{The canonical volume and characteristic pairings.}
On the torus cover away from $q=0$ consider
\[
                    \omega=d\log X\wedge d\log Y\wedge dq.
                                                               \tag{D47.10}
\]
The period translations multiply $X,Y$ by integral
powers of $q$ and by nonvanishing functions of $q$.
The added logarithmic differential terms are multiples
of $dq$, so $\omega$ is invariant and descends.
In every unimodular triangular chart one can choose
coordinates with $q=xyz$, $X=x$, $Y=y$ up to the
unimodular monomial change. Then
\[
                    d\log x\wedge d\log y\wedge d(xyz)
                                 =dx\wedge dy\wedge dz.
                                                               \tag{D47.11}
\]
Thus the descended form extends without a zero or pole
over the central fibre, including self-intersections.
The cap has a trivial canonical bundle near all its
compact components.

The normalization $\nu_i$ is a holomorphic immersion:
each preimage of a crossing lies on one smooth local
coordinate branch. Even when its image has self-gluings,
the immersion normal line is therefore
$\mathcal N_{\nu_i}=K_{P_i}$, by the adjunction sequence
and the nonvanishing volume form. This normal line
must not be confused with $\nu_i^*\mathcal O(E_i)$,
which includes the other self-intersection branches.
The exact sequence
\[
       0\longrightarrow TP_i\longrightarrow\nu_i^*TN_L
                        \longrightarrow K_{P_i}\longrightarrow0
                                                               \tag{D47.12}
\]
now gives
\[
 \nu_i^*c_2(TN_L)=c_2(TP_i)-c_1(TP_i)^2,\qquad
 \nu_i^*p_1(TN_L)=2c_1(TP_i)^2-2c_2(TP_i).                 \tag{D47.13}
\]
The hexagonal surface has Euler number six (one for
each torus-fixed point) and $c_1^2=6$, also obtained
by squaring $\sum_rC_{i,r}$ using (D47.6).
Both evaluations in (D47.13) are therefore zero,
proving (D47.4). The calculation uses the normal line
of the immersion and is valid for every quotient
lattice, not only for embedded central components.

\emph{Exact verification and bounds.}
The program
\path{code/fourfolds_D2_component_ring.py}
enumerates all $127$ column Hermite forms
$\left(\begin{smallmatrix}a&b\\0&d\end{smallmatrix}\right)$
with $ad\le12$ and $0\le b<a$. For every one it
computes the integral triple products directly from
the six boundary restrictions, verifies symmetry
in all three indices and the constant annihilator,
and compares every cubic coefficient with the
quotient cell formula while retaining loops and
repeated labels. It also verifies the toric ray,
intersection and canonical-volume identities and
the examples (D47.5). All matrices, cells and
nonzero integral products are retained in
\path{records/fourfolds_D2_component_ring.json}.
The proof above applies to arbitrary finite index;
the finite computation tests the conventions and
does not supply an unstated global index bound.
\hfill$\square$
